\documentclass[11pt,a4paper]{article}
\usepackage[T1]{fontenc}
\usepackage[utf8]{inputenc}
\usepackage{amsmath,amssymb,amsthm}
\usepackage[margin=1in]{geometry}
\usepackage[colorlinks=true,linkcolor=blue,urlcolor=blue]{hyperref}
\theoremstyle{plain}
\newtheorem{theorem}{Theorem}
\newtheorem{lemma}[theorem]{Lemma}
\newtheorem{proposition}[theorem]{Proposition}
\newtheorem{corollary}[theorem]{Corollary}
\theoremstyle{definition}
\newtheorem{remark}[theorem]{Remark}

\title{The empirical recurrence for OEIS A209099, proved by counting patterns}
\author{Adrian Perez Fontelles\\ \small Independent researcher}
\date{31 August 2026}
\begin{document}
\maketitle

\begin{abstract}
OEIS A209099 counts arrays subject to a local condition \emph{and} to the clause ``new values
$0..2$ introduced in row major order'', which says the array is written in canonical
form. That clause is not local, so a transfer matrix over the alphabet does not see it. It
does not have to: what is being counted is not arrays but equality patterns, and the number
of patterns is recovered from the numbers $L_i$ of arrays over an $i$-letter alphabet by
inverting a falling factorial. The result is the clean identity
$3!\,a(n)=\sum_i\binom{3}{i}D_{3-i}L_i(n)$ with $D_m$ the derangement numbers,
and each $L_i$ is an ordinary walk count. A second reduction --- the chain is strongly
lumpable over the patterns, because the condition is relabelling-invariant --- brings the
state space down from $2316$ to $S=430$. The entry's empirical recurrence of order
$62$, contributed by R.~H.~Hardin, then follows from a finite exact computation.
\end{abstract}

\noindent\small 2020 Mathematics Subject Classification. 05A15, 05A18, 15B36.\normalsize

\section{The sequence and the conjecture}

OEIS A209099 is ``Number of nX7 0..2 arrays with new values 0..2 introduced in row major order and no element equal to more than one of its immediate leftward or upward or right-upward antidiagonal neighbors.''. It has offset $1$ and begins
\[
365, 165632, 81254376, 40556341276, 20321585350224, 10191444894367900, 5112092569118325704, 2564363579403254284620,\ \dots
\]
The entry carries the following comment:
\begin{quote}\small
Empirical: a(n) = 747*a(n-1) -147770*a(n-2) +13767161*a(n-3) -744653684*a(n-4) +26581250215*a(n-5) -695207333006*a(n-6) +14759817640723*a(n-7) -277760418691959*a(n-8) +4818498687693894*a(n-9) -76394118399894709*a(n-10) +1087272633889124267*a(n-11) -13781949249633025934*a(n-12) +154877879063702829937*a(n-13) -1530486551719044001957*a(n-14) +13153680955142056537404*a(n-15) -97057933858981232648357*a(n-16) +604009502946148839126239*a(n-17) -3089743014121715957566207*a(n-18) +12592151665916950689366304*a(n-19) -39754681120743615063631342*a(n-20) +95137341741302104796392420*a(n-21) -157046844212176575473908358*a(n-22) +62809971153790464358833542*a(n-23) +274260546953763747783830412*a(n-24) +2277040460294389799199241820*a(n-25) -16807391841560612463088548960*a(n-26) +24466406030744021830227782360*a(n-27) +110770739770822017548456685488*a(n-28) -642543090446277892020887588984*a(n-29) +1444810249046399642357458572124*a(n-30) -1039066917861135086134740251232*a(n-31) -1913422528148454603288634391464*a(n-32) +4382959090538650287480505632320*a(n-33) -420162726254070213081625756976*a(n-34) -10248211114440498997980265577120*a(n-35) +19099890867084707788818392405696*a(n-36) -18080871254247582043431141300352*a(n-37) +9535746814747597050087649329920*a(n-38) -7274842106933434445259296042752*a(n-39) +24412291656612509021701274945408*a(n-40) -50569037378415505084329421866496*a(n-41) +49696873654792766125531687784960*a(n-42) -3207612319613312961503561953792*a(n-43) -54507430486119902763301778119680*a(n-44) +71876543544959492841651732609024*a(n-45) -41832833113226911264957080469504*a(n-46) +2042228750173828354218287046656*a(n-47) +16335249204577141685269880815616*a(n-48) -14493780574353894230229010939904*a(n-49) +7212093727799570963008951418880*a(n-50) -2258955815739930161535614189568*a(n-51) +173106370971311909976634556416*a(n-52) +236456074127730455885312425984*a(n-53) -132988789131444826419830456320*a(n-54) +57804329724143223598793883648*a(n-55) -21600473043383892821286060032*a(n-56) -190229647086921884471656448*a(n-57) +3407643142423947053302284288*a(n-58) -1076075759469831874666299392*a(n-59) +124348221639171208497856512*a(n-60) -2055059748711593966305280*a(n-61) -370952494107253014528000*a(n-62) for n>65
\end{quote}
As of the ``Last modified'' line on the live entry (Oct 03 2025 15:09:03, revision 9) the statement is
still recorded as empirical, and nothing on the entry records it as settled either way.

Written out, the claim is
\begin{equation}\label{eq:conj}
a(n)\;=\;747\,a(n-1) - 147770\,a(n-2) + 13767161\,a(n-3) - 744653684\,a(n-4) + 26581250215\,a(n-5) - 695207333006\,a(n-6) + 14759817640723\,a(n-7) - 277760418691959\,a(n-8) + 4818498687693894\,a(n-9) - 76394118399894709\,a(n-10) + 1087272633889124267\,a(n-11) - 13781949249633025934\,a(n-12) + 154877879063702829937\,a(n-13) - 1530486551719044001957\,a(n-14) + 13153680955142056537404\,a(n-15) - 97057933858981232648357\,a(n-16) + 604009502946148839126239\,a(n-17) - 3089743014121715957566207\,a(n-18) + 12592151665916950689366304\,a(n-19) - 39754681120743615063631342\,a(n-20) + 95137341741302104796392420\,a(n-21) - 157046844212176575473908358\,a(n-22) + 62809971153790464358833542\,a(n-23) + 274260546953763747783830412\,a(n-24) + 2277040460294389799199241820\,a(n-25) - 16807391841560612463088548960\,a(n-26) + 24466406030744021830227782360\,a(n-27) + 110770739770822017548456685488\,a(n-28) - 642543090446277892020887588984\,a(n-29) + 1444810249046399642357458572124\,a(n-30) - 1039066917861135086134740251232\,a(n-31) - 1913422528148454603288634391464\,a(n-32) + 4382959090538650287480505632320\,a(n-33) - 420162726254070213081625756976\,a(n-34) - 10248211114440498997980265577120\,a(n-35) + 19099890867084707788818392405696\,a(n-36) - 18080871254247582043431141300352\,a(n-37) + 9535746814747597050087649329920\,a(n-38) - 7274842106933434445259296042752\,a(n-39) + 24412291656612509021701274945408\,a(n-40) - 50569037378415505084329421866496\,a(n-41) + 49696873654792766125531687784960\,a(n-42) - 3207612319613312961503561953792\,a(n-43) - 54507430486119902763301778119680\,a(n-44) + 71876543544959492841651732609024\,a(n-45) - 41832833113226911264957080469504\,a(n-46) + 2042228750173828354218287046656\,a(n-47) + 16335249204577141685269880815616\,a(n-48) - 14493780574353894230229010939904\,a(n-49) + 7212093727799570963008951418880\,a(n-50) - 2258955815739930161535614189568\,a(n-51) + 173106370971311909976634556416\,a(n-52) + 236456074127730455885312425984\,a(n-53) - 132988789131444826419830456320\,a(n-54) + 57804329724143223598793883648\,a(n-55) - 21600473043383892821286060032\,a(n-56) - 190229647086921884471656448\,a(n-57) + 3407643142423947053302284288\,a(n-58) - 1076075759469831874666299392\,a(n-59) + 124348221639171208497856512\,a(n-60) - 2055059748711593966305280\,a(n-61) - 370952494107253014528000\,a(n-62)
\end{equation}
for all $n>65$.

\section{What the canonical-form clause counts}

Fix the shape: $7$ columns and $L=n+0$ rows, so $L\cdot7$ cells. Call a
partition of the cells into nonempty classes a \emph{pattern}; an array over an alphabet
determines a pattern, two cells lying in the same class exactly when they carry the same
value.

Write $x_{t,u}$ for the entry in row $t$ and column $u$. The
entry's requirement is that no cell equals more than $1$ of its neighbours at the
offsets
\[
\mathcal N=\{(0,-1),\ (-1,0),\ (-1,1)\}
\]
(offsets given as (row shift, column shift), and only cells inside the array count),
that is
\begin{equation}\label{eq:loc}
\#\bigl\{(d_1,d_2)\in\mathcal N:\ x_{t+d_1,\,u+d_2}=x_{t,u}\bigr\}\;\le\;1
\qquad\text{for every cell }(t,u).
\end{equation}
Every offset in $\mathcal N$ points backwards in the walk, so the condition at a cell of
row $t$ is settled by row $t$ itself together with row $t-1$.

The condition is stated purely in terms of equalities between entries, so it is invariant
under any relabelling of the values and is therefore a property of the pattern alone. Call a
pattern \emph{admissible} when it has the property.

The clause ``new values $0..2$ introduced in row major order'' selects, from each class
of arrays that differ only by a relabelling, exactly one representative: reading the cells in
row major order, the first occurrence of the value $v$ must precede the first occurrence of
$v+1$. Every pattern with at most $K=3$ classes has exactly one such representative, and
every array in canonical form is the representative of its own pattern. Hence, writing
$N_j(n)$ for the number of admissible patterns with exactly $j$ classes,
\begin{equation}\label{eq:apat}
a(n)\;=\;\sum_{j=0}^{3}N_j(n).
\end{equation}

\begin{lemma}\label{lem:inv}
Let $L_i(n)$ be the number of arrays over an alphabet of $i$ letters satisfying the local
condition, with no canonical-form clause. Then
\[
L_i(n)=\sum_j N_j(n)\,i(i-1)\cdots(i-j+1),
\qquad
3!\,\sum_{j=0}^{3}N_j(n)=\sum_{i=0}^{3}\binom{3}{i}D_{3-i}\,L_i(n),
\]
where $D_m=m!\sum_{t=0}^{m}(-1)^t/t!$ is the number of derangements of $m$ letters.
\end{lemma}

\begin{proof}
An array over $i$ letters satisfying the condition determines an admissible pattern together
with an injection of its classes into the alphabet, and conversely; a pattern with $j$
classes admits $i(i-1)\cdots(i-j+1)$ injections. That is the first identity. Writing
$i(i-1)\cdots(i-j+1)=j!\binom{i}{j}$ it reads $L_i=\sum_j (j!N_j)\binom{i}{j}$, so by
binomial inversion $j!N_j=\sum_{i}(-1)^{j-i}\binom{j}{i}L_i$. Summing over
$j\le 3$ and exchanging the order of summation,
\[
\sum_{j=0}^{3}N_j=\sum_{i=0}^{3}L_i\sum_{j=i}^{3}\frac{(-1)^{j-i}}{i!\,(j-i)!}
=\sum_{i=0}^{3}\frac{L_i}{i!}\sum_{t=0}^{3-i}\frac{(-1)^t}{t!} ,
\]
and multiplying by $3!$ turns the coefficient of $L_i$ into
$\frac{3!}{i!}\sum_{t\le 3-i}(-1)^t/t!=\binom{3}{i}(3-i)!\sum_{t\le 3-i}(-1)^t/t!
=\binom{3}{i}D_{3-i}$, an integer.
\end{proof}

\section{Each $L_i$ is a walk count, and the chain lumps}

Fix $i$ and let $V_i=\{0,\dots,i-1\}^{7}$ be the possible rows. For $r,s\in V_i$
declare $r\to s$ an edge when placing $s$ directly after $r$ satisfies \eqref{eq:loc} at
every column, and let $M_i$ be the adjacency matrix. An array is exactly a sequence of
rows, and its condition is exactly the conjunction of the edge conditions along that
sequence, so
\[
L_i(n)\;=\;\mathbf{s}_i^{\!\top}M_i^{\,n+0-1}\mathbf{1} .
\]
The first row has no row before it, so its cells are tested
against their within-row neighbours only; $\mathbf{s}_i$ is the indicator vector of the
rows that pass that test.

\begin{lemma}\label{lem:lump}
Let $\pi(r)$ be the equality pattern of the row $r$. For all $r,r'$ with
$\pi(r)=\pi(r')$ and every pattern $Q$,
\[
\#\{s:\ r\to s,\ \pi(s)=Q\}\;=\;\#\{s:\ r'\to s,\ \pi(s)=Q\}\, .
\]
Consequently the numbers $\widehat M_i[P][Q]$ obtained by taking any representative of $P$
are well defined, and with $e_P=i(i-1)\cdots(i-|P|+1)$ the number of rows of pattern $P$,
\[
\mathbf{s}_i^{\!\top}M_i^{\,m}\mathbf{1}\;=\;\sum_P \varepsilon_P\,f_m(P),
\qquad f_0\equiv1,\quad f_{m+1}(P)=\sum_Q\widehat M_i[P][Q]f_m(Q),
\]
where $\varepsilon_P=e_P$ if rows of pattern $P$ may start a walk and $0$ otherwise.
\end{lemma}

\begin{proof}
If $\pi(r)=\pi(r')$ then the map sending the value of $r$ at a column to the value of $r'$
at that column is a well-defined injection between the letters used, and extends to a
permutation $\sigma$ of the alphabet with $r'=\sigma(r)$. The edge condition is a statement
about equalities among entries, so $r\to s$ if and only if $\sigma(r)\to\sigma(s)$; and
$\pi(\sigma(s))=\pi(s)$. Hence $s\mapsto\sigma(s)$ is a bijection between the two successor
sets preserving patterns, which is the first claim. It says the chain is strongly lumpable
over $\pi$, so $\bigl(M_i^{m}\mathbf{1}\bigr)(r)$ depends on $r$ only through $\pi(r)$;
calling that value $f_m(\pi(r))$ gives the recursion, and summing over $r$ in each class
gives the displayed identity, the number of rows of pattern $P$ being the number of
injections of its $|P|$ classes into the alphabet.
\end{proof}

Assembling the $3$ lumped chains into one block-diagonal matrix $N$ of size $S=430$ and
putting the weights of Lemma~\ref{lem:inv} into the start vector,
$\mathbf{v}=\bigl(\binom{3}{i}D_{3-i}\varepsilon_P\bigr)_{i,P}$, gives
\begin{equation}\label{eq:walk}
3!\;a(n)\;=\;\mathbf{v}^{\!\top}N^{\,n+0-1}\mathbf{1} .
\end{equation}
The unlumped alphabet had $2316$ rows in all; the lumped chain has $S=430$ states.

\section{The criterion, and the computation}

Let $q(t)=t^{62}-\sum_i c_i\,t^{62-i}$ be the characteristic polynomial of
\eqref{eq:conj}.

\begin{lemma}\label{lem:crit}
For every $n$ with $m=n+0-1\ge62$,
\[
3!\Bigl(a(n)-\sum_i c_i\,a(n-i)\Bigr)\;=\;\mathbf{v}^{\!\top}N^{\,m-62}q(N)\mathbf{1} .
\]
Hence \eqref{eq:conj} holds from that point on if and only if
$\mathbf{v}^{\!\top}N^{j}q(N)\mathbf{1}=0$ for every $j\ge0$, and it is enough to check
$j=0,1,\dots,S-1$.
\end{lemma}

\begin{proof}
Substitute \eqref{eq:walk} into the left side and factor $N^{m-62}$ out of
$N^{m}-\sum_i c_iN^{m-i}$. The scalar $3!$ is nonzero, so it does not affect whether
the left side vanishes. For the last clause put $w=q(N)\mathbf{1}$: by Cayley--Hamilton
$N^{S}$ is an integer combination of $I,N,\dots,N^{S-1}$, so each
$\mathbf{v}^{\!\top}N^{j}w$ with $j\ge S$ is the corresponding combination of earlier
values.
\end{proof}

\begin{theorem}
The recurrence \eqref{eq:conj} holds for every $n>65$, which is the statement of
Section~1.
\end{theorem}

\begin{proof}
The lumped transition counts were obtained by enumerating, for one representative row of
each pattern, all admissible successors and sorting them by pattern --- the successors being
generated column by column, since the edge condition is a conjunction of one constraint per
column. The vector $w=q(N)\mathbf{1}$ was then formed by $62$ matrix--vector products
in exact integer arithmetic and $\mathbf{v}^{\!\top}N^{j}w$ evaluated for
$j=0,1,\dots$, again exactly. Every one of those integers vanishes from the index
corresponding to $n=65+1$ onwards, and the run continued until $w$ was the zero vector,
which settles all larger $j$ at once. Lemma~\ref{lem:crit} gives the theorem.
\end{proof}

\section{Verification}

Three checks, each independent of the algebra above.

First, the whole construction was compared against the entry's DATA: the right-hand side of
\eqref{eq:walk}, divided by $3!$, was evaluated at every one of the $11$
published indices and agrees exactly. This is a strong test of the modelling and not only of
the arithmetic --- the inversion of Lemma~\ref{lem:inv}, the reading of the local
condition, and the alignment of the entry's index with the number of rows would each
show up as a mismatch, and the quotient by $3!$ came out an integer at every index.

Second, the lumped chain was checked against the unlumped one: for the small shapes the walk
counts $\mathbf{s}_i^{\!\top}M_i^{m}\mathbf{1}$ were computed directly over the full
alphabet as well and agree with $\sum_P\varepsilon_Pf_m(P)$ term by term.

Third, the recurrence \eqref{eq:conj} was evaluated on the published terms in exact integer
arithmetic with no matrices involved, and the annihilation test of Lemma~\ref{lem:crit} was
run until the working vector was identically zero rather than to a sampled prefix.

\begin{thebibliography}{9}
\bibitem{oeis} The OEIS Foundation, \emph{The On-Line Encyclopedia of Integer
Sequences}, \url{https://oeis.org}, 2026. Sequence A209099.
\bibitem{stanley} R.~P.~Stanley, \emph{Enumerative Combinatorics}, Volume 1, 2nd ed.,
Cambridge University Press, 2011. (Transfer-matrix method, Section 4.7; Stirling and
falling-factorial inversion, Section 1.9.)
\bibitem{kemeny} J.~G.~Kemeny and J.~L.~Snell, \emph{Finite Markov Chains}, Springer,
1976. (Strong lumpability, Section 6.3.)
\bibitem{fs} P.~Flajolet and R.~Sedgewick, \emph{Analytic Combinatorics}, Cambridge
University Press, 2009.
\end{thebibliography}

\end{document}
